\documentclass{article}
\usepackage[utf8]{inputenc}
\usepackage{amsmath}
\usepackage{graphicx}
\usepackage{subcaption}
\usepackage{float}
\usepackage{geometry}
\geometry{a4paper, margin=1in}

\title{SH1017 Physics for Engineering Mathematics \\ Computer Lab 2: Diffraction and Interference}
\author{--}
\date{\today}

\begin{document}

\maketitle

\section*{Background}
This report details the findings from Computer Lab 2, focusing on the simulation of wave interference and diffraction using point sources. We investigate the behavior of interference patterns as we vary source configurations, slit widths, and distances. The electric field $E(\mathbf{r}, t)$ from $N$ identical sources at positions $\mathbf{r}_i$ is given by the superposition:
\begin{equation}
E(\mathbf{r}, t) = \sum_{i=1}^{N} \frac{A}{|\mathbf{r} - \mathbf{r}_i|} \cos(k|\mathbf{r} - \mathbf{r}_i| - \omega t)
\end{equation}
The time-averaged intensity measured on a detector screen at distance $D$ is:
\begin{equation}
\overline{E^2} = \frac{|A|^2}{r^2} \left( \frac{N}{2} + \sum_{i=1}^{N-1} \sum_{j=i+1}^N \cos(k|\mathbf{r} - \mathbf{r}_i| - k|\mathbf{r} - \mathbf{r}_j|) \right)
\end{equation}

\section*{Q1: Interference from Point Sources}

\textbf{Method:} Using \texttt{interference.py}, we simulated two coherent point sources and observed how the interference pattern changes with the separation distance $d$ and by adding a third source.

\begin{figure}[H]
    \centering
    \begin{subfigure}[b]{0.45\textwidth}
        \includegraphics[width=\textwidth]{figures/p1_d10.png}
        \caption{$d=5$}
    \end{subfigure}
    \begin{subfigure}[b]{0.45\textwidth}
        \includegraphics[width=\textwidth]{figures/p1_d20.png}
        \caption{$d=25$}
    \end{subfigure}
    \begin{subfigure}[b]{0.45\textwidth}
        \includegraphics[width=\textwidth]{figures/p1_d40.png}
        \caption{$d=50$}
    \end{subfigure}
    \begin{subfigure}[b]{0.45\textwidth}
        \includegraphics[width=\textwidth]{figures/p1_3sources.png}
        \caption{3 sources, $d=25$}
    \end{subfigure}
    \caption{Interference patterns for different configurations ($\lambda=5.0$ mm, $100 \times 100$ mm area).}
    \label{fig:p1}
\end{figure}

\textbf{Theory:} According to $N$-slit interference theory, the condition for constructive interference (maxima) is given by $d \sin \theta = m \lambda$, where $m$ is an integer. Thus, increasing the distance $d$ between sources must decrease the angular separation $\theta$ between fringes, leading to more fringes within a given area. Furthermore, the intensity for $N$ sources scales as $I \propto \frac{\sin^2(N \delta/2)}{\sin^2(\delta/2)}$ where $\delta = kd \sin\theta$. Adding more sources (such as a third source, $N=3$) therefore sharpens the primary maxima and introduces $(N-2)$ secondary, smaller maxima between them.

\textbf{Results:} As seen in Figure \ref{fig:p1}, increasing the distance $d$ between the sources increases the number of interference fringes (maxima and minima) within the same area, corresponding to a smaller angular separation. Adding a third source between the two existing ones sharpens the primary maxima and introduces secondary, smaller maxima between them. This is perfectly consistent with theoretical predictions.

\section*{Q2: Double Slit Interference}

\textbf{Method:} Using the code from \texttt{doubleslit.py} we simulated the interference pattern on a screen at distance $D$ from two slits separated by $d$, using $\lambda = 500$ nm. We compared the measured fringe spacing $\Delta x$ with the theoretical result $\Delta x = \lambda D / d$.

\textbf{Theory:} The condition for maxima is path difference $\Delta r \approx d \sin \theta = m \lambda$. Using a small-angle approximation ($\sin \theta \approx \tan \theta = x/D$), we derive the theoretical fringe spacing formula:
\begin{equation*}
    \Delta x = \frac{\lambda D}{d}
\end{equation*}
This formula implies a linear relationship between the fringe position and interference order on a flat screen. However, without these approximations, the exact position of maxima on a flat screen is $x_m = D \tan(\arcsin(m\lambda/d))$. Because $\tan \theta$ grows faster than $\theta$, the distance between fringes on a flat screen should increase as we move further from the center compared to what the simple linear formula predicts.

\textbf{Results:} The results for various parameters are summarized below:
\begin{table}[h]
\centering
\begin{tabular}{|l|c|c|c|c|c|}
\hline
Case & $\lambda$ (mm) & $d$ (mm) & $D$ (mm) & Theoretical $\Delta x$ (mm) & Measured $\Delta x$ (mm) \\ \hline
Base & 0.0005 & 0.5 & 2000 & 2.0000 & 2.6627 \\
Red light & 0.0007 & 0.5 & 2000 & 2.8000 & 5.5856 \\
Small $d$ & 0.0005 & 0.25 & 2000 & 4.0000 & 7.9880 \\
Small $D$ & 0.0005 & 0.5 & 1000 & 1.0000 & 1.1100 \\
\hline
\end{tabular}
\caption{Comparison of theoretical and measured fringe spacing.}
\label{tab:p2}
\end{table}

The measured values differ and are consistently larger than the simple formula $\Delta x = \lambda D / d$. This is because the formula assumes a small-angle approximation ($\sin \theta \approx \theta$) and a planar detector screen, whereas the simulation uses the exact geometric distance $r = \sqrt{D^2 + x^2}$. As expected, as $x$ increases, the distance between fringes on a flat screen increases compared to an arc, fully explaining the observed deviations from the theoretical approximation.

\section*{Q3: Single Slit Model}

\textbf{Method:} We modified \texttt{doubleslit.py} to approximate a single slit of width $w = 0.05$ mm by placing $N \in \{2, 3, 5, 10\}$ point sources along the slit, using $\lambda = 0.0005$ mm and $D=2000$ mm. We investigated how the peak intensity and width depend on $N$ and $w$.

\begin{figure}[H]
    \centering
    \begin{subfigure}[b]{0.45\textwidth}
        \includegraphics[width=\textwidth]{figures/p3_N2.png}
        \caption{$N=2$}
    \end{subfigure}
    \begin{subfigure}[b]{0.45\textwidth}
        \includegraphics[width=\textwidth]{figures/p3_N3.png}
        \caption{$N=3$}
    \end{subfigure}
    \begin{subfigure}[b]{0.45\textwidth}
        \includegraphics[width=\textwidth]{figures/p3_N5.png}
        \caption{$N=5$}
    \end{subfigure}
    \begin{subfigure}[b]{0.45\textwidth}
        \includegraphics[width=\textwidth]{figures/p3_N10.png}
        \caption{$N=10$}
    \end{subfigure}
    \caption{Intensity distribution for a single slit modeled with varying $N$ sources ($w=0.05$ mm, $\lambda=0.0005$ mm, $D=2000$ mm).}
    \label{fig:p3_N}
\end{figure}

\begin{figure}[H]
    \centering
    \begin{subfigure}[b]{0.45\textwidth}
        \includegraphics[width=\textwidth]{figures/p3_N_diff_w2.png}
        \caption{$N=2$}
    \end{subfigure}
    \begin{subfigure}[b]{0.45\textwidth}
        \includegraphics[width=\textwidth]{figures/p3_N_diff_w3.png}
        \caption{$N=3$}
    \end{subfigure}
    \begin{subfigure}[b]{0.45\textwidth}
        \includegraphics[width=\textwidth]{figures/p3_N_diff_w5.png}
        \caption{$N=5$}
    \end{subfigure}
    \begin{subfigure}[b]{0.45\textwidth}
        \includegraphics[width=\textwidth]{figures/p3_N_diff_w10.png}
        \caption{$N=10$}
    \end{subfigure}
    \caption{Intensity distribution for a single slit modeled with varying $N$ sources ($w=0.2$ mm, $\lambda=0.0005$ mm, $D=2000$ mm).}
    \label{fig:p3_w}
\end{figure}

\textbf{Theory:} Adjusting the number of point sources $N$ inside a fixed slit width $w$ determines how well we approximate a continuous slit. The peak intensity at the center ($x=0$) scales quadratically with the number of sources:
\begin{equation*}
    I_{max} \propto \left(\sum_{i=1}^N A\right)^2 = N^2 A^2
\end{equation*}
because the $N$ coherent waves perfectly constructively interfere. As $N \to \infty$, the discrete sum converges to the continuous single-slit diffraction intensity formula:
\begin{equation*}
    I(\theta) = I_0 \operatorname{sinc}^2\left(\frac{\pi w \sin \theta}{\lambda}\right)
\end{equation*}
This continuous pattern exhibits distinct side lobes. Furthermore, the first minima occur at $w \sin \theta = \pm \lambda$, meaning the angular width of the central diffraction peak is inversely proportional to the slit width $w$ (specifically, the angular half-width is $\theta \approx \lambda/w$). Thus, widening the slit $w$ will result in a narrower central peak, while narrowing $w$ will yield a wider central peak.

\textbf{Results:} As seen in Figure \ref{fig:p3_N} ($w=0.05$ mm), increasing $N$ correctly causes the central peak intensity to grow quadratically, matching the expected $N^2$ scaling behavior (for example, the peak for $N=10$ is exactly 25 times taller than the peak for $N=2$). Moreover, as $N$ increases to 10, we see the pattern clearly evolve to exhibit secondary side lobes, approaching the true single-slit diffraction envelope. 

Figure \ref{fig:p3_w} demonstrates the same variation in $N$ but for a wider slit ($w=0.2$ mm). Comparing the fully developed patterns ($N=10$) between Figure \ref{fig:p3_w} and Figure \ref{fig:p3_N}, the central peak for $w=0.2$ mm is visibly much narrower than for $w=0.05$ mm, and its side lobes are packed more tightly together. These observations demonstrate that the peak width is indeed inversely proportional to the slit width $w$, which strongly confirms both the scaling and interference principles of single-slit diffraction theory.

\section*{Q4: Realistic Double Slit Model}

\textbf{Method:} We copied 2 of the "realistic" slits from problem 3 ($N=10$, $w=0.05$ mm) and separated them by $d=0.5$ mm and compared it to the point-source model.

\begin{figure}[H]
    \centering
        \begin{subfigure}[b]{0.8\textwidth}
        \includegraphics[width=\textwidth]{figures/p2_base.png}
        \caption{Point source model}
    \end{subfigure}
    \begin{subfigure}[b]{0.8\textwidth}
        \includegraphics[width=\textwidth]{figures/p4_comparison.png}
        \caption{Realistic slit model}
    \end{subfigure}
    \caption{Comparison between point source double slit ($d=0.5$ mm, $\lambda=0.0005$ mm, $D=2000$ mm) and realistic double slit (with discretization along a continuous width $N=10$, $w=0.05$ mm).}
    \label{fig:p4}
\end{figure}

\textbf{Theory:} A realistic double slit model should incorporate both interference and diffraction effects. Theoretical principles suggest that the ideal interference fringes will be modulated by the single-slit diffraction envelope, yielding a combined intensity formula:
\begin{equation*}
    I(\theta) = I_0 \underbrace{\cos^2\left(\frac{\pi d \sin \theta}{\lambda}\right)}_{\text{interference}} \underbrace{\operatorname{sinc}^2\left(\frac{\pi w \sin \theta}{\lambda}\right)}_{\text{diffraction}}
\end{equation*}
Consequently, the interference fringes (from the $\cos^2$ term) should diminish in intensity further from the central peak, disappearing entirely at the diffraction envelope's null points (where the $\operatorname{sinc}^2$ term goes to zero).

\textbf{Results:} As shown in Figure \ref{fig:p4}, subfigure (a) depicts the ideal double slit modeled with perfect point sources, which produces interference fringes of uniform maximum intensity across the entire screen. Subfigure (b) depicts the realistic slit model, which introduces the finite width $w=0.05$ mm. Comparing the two, the realistic model maintains the exact same interference fringe spacing ($\Delta x$) as the ideal model. However, the intensity of these fringes is now correctly bounded and modulated by the broader single-slit diffraction envelope. This causes the individual fringes to progressively decrease in intensity as they move away from the central maximum, eventually disappearing entirely at the nulls of the diffraction envelope. This behavior demonstrates the physical convolution of interference and diffraction, in full agreement with the realistic theoretical model.

\end{document}
